% ============================================================
\chapter*{この資料について}
\addcontentsline{toc}{chapter}{この資料について}

Black-Litterman モデル（以下 BL）は、市場の時価総額の比率から逆算した「均衡リターン」を出発点にして、投資家の見方（ビュー）を確信の度合いに応じて混ぜ、期待リターンを作り直す方法である。式そのものは短い。難しいのは、式に入れる値を決めることである。とくに次の 3 つは、論文でも実務でも決め方が定まっていない。

\begin{itemize}
\item \textbf{ビューの値 $Q$}：何をどれだけ上がる（下がる）と見るか。
\item \textbf{ビューの不確かさ $\Omega$}：そのビューをどれだけ信じるか。
\item \textbf{均衡の不確かさの大きさ $\tau$}：出発点の均衡リターンをどれだけ信じるか。
\end{itemize}

この資料は、この 3 つ（と、それに関わる共分散 $\Sigma$、ビューの形 $P$）の決め方の「コツ」を、論文を読んで整理し、本プロジェクト（LLM-BLM）のデータで実際に計算して確かめるものである。

\section*{ほかの資料との関係}

\begin{center}
\small
\begin{tabular}{p{3.0cm}p{10.6cm}}
\toprule
資料 & 内容 \\
\midrule
『BLM解説』（勉強用） & 平均分散最適化から BL までの基礎。He と Litterman（1999）に沿って、公式の導き方と G7 の例を手で計算する \\
『総合報告書』 & 対象論文（arXiv 2504.14345）の再現と検証、データ基盤、実験 A〜G。元の計算の問題点（ビューが 99.9\% 採用される、回答のばらつきは信頼度ではない、など） \\
この資料 & ビュー・不確かさ・$\tau$ の決め方。論文ごとの方法の整理と、その数値例。本プロジェクトの LLM の回答で $\Omega$ の決め方を比べ直す \\
\bottomrule
\end{tabular}
\end{center}

基礎は『BLM解説』と重なるところもあるが、この資料だけで読めるように、必要な式は導き直している。

\section*{依頼のときに渡された調査メモについて}

この資料は、依頼のときに渡された文章（以下「調査メモ」）を骨組みにしている。調査メモは、BL のビューとリスクを数値化する方法を、政治・マクロ経済・財務（ESG を含む）・モメンタムの 4 つのファクターに分けて整理し、ファクター版の BL、エントロピー・プーリング、トラッキングエラー付きの最適化にも触れている。

調査メモには、確かめられる主張と、確かめられない主張が混ざっている。この資料では、論文の本文または要旨で確認できたことだけを事実として書いた。確認できなかったもの、論文と食い違うものは、\appref{app:check} に一覧にした。主なものを挙げる。

\begin{itemize}
\item Sun ほか（2026）の「年率 106.58\%、シャープレシオ 2.64」は論文の数字どおりである。ただし、対象は 8 銘柄、評価期間は 117 取引日で、ビューの不確かさは単位行列と書かれている（\chapref{ch:text}）。
\item 対象論文が回答を 100 回取るのは第 2 版である。第 1 版と本プロジェクトは 30 回である。
\item 「ESG の点数で共分散を $(2-\mathrm{ESG})$ 倍に膨らませる」方法と、「景気循環のビューと CVaR を組み合わせてシャープレシオが平均 19.7\% 改善した」という研究は、出典を見つけられなかった（\chapref{ch:esg}、\chapref{ch:macro}）。
\item 「LLM の回答のばらつきが、情報の不確かさを客観的に表す」という主張は、本プロジェクトのデータでは成り立たない。回答の分散は、実際の予想の誤差の 2 乗の中央値で \V{llm.ratio.med} 倍しかない（\chapref{ch:project}）。
\end{itemize}

\section*{構成}

\begin{center}
\small
\begin{tabular}{p{2.4cm}p{11.2cm}}
\toprule
部 & 内容 \\
\midrule
第 I 部 基礎 & BL の全体像。公式の 3 通りの導き方。均衡リターンとリスク回避度。$\tau$ の意味と決め方 \\
第 II 部 $\Omega$ の決め方 & 方法の一覧。Idzorek の確信度の方法を例で再現する。He--Litterman と Meucci の流儀。予想の当たり具合（IC）から決める方法。ビューが多いときの確信の積み上がり。ビューの形（絶対・相対・順位・ファクター） \\
第 III 部 ファクターごとの数値化 & 政治と政策の不確かさ（GPR、EPU）。テキストのセンチメントと LLM。マクロとボラティリティ（EGARCH-M）。財務と ESG。モメンタム \\
第 IV 部 拡張と最適化 & エントロピー・プーリング。コピュラ・オピニオン・プーリング。CVaR の最適化。トラッキングエラーと情報レシオ \\
第 V 部 当てはめ & 本プロジェクトの LLM のビューで $\Omega$ の決め方を比べる。コツの一覧。まとめ \\
付録 & 行列の公式。調査メモの主張の確認表。論文の一覧。作り直し方。Python での実装例。用語集。参考文献 \\
\bottomrule
\end{tabular}
\end{center}

急ぐ読者は、\chapref{ch:overview}（全体像）、\chapref{ch:accuracy}（当たり具合から $\Omega$ を決める）、\chapref{ch:project}（本プロジェクトでの比較）、\chapref{ch:tips}（コツの一覧）を読めば要点がつかめる。

\section*{数字の出どころ}

本文の数値・表・図は、すべて \code{ビューとリスクのコツ\_計算.py} が計算して \code{生成/} に書き出したものである。本文に数字を手で書き写していない。使ったデータは次のとおりである（詳しくは \appref{app:rebuild}）。

\begin{itemize}
\item 本プロジェクトの DB（\code{data/blm\_2015.sqlite}）：株価、指数、セクター ETF、\SP{} の構成銘柄と時価総額、LLM の回答（実行番号 1、gpt-4o-mini、30 か月、約 44 万個）。
\item Kenneth French のデータライブラリ：Fama--French の 5 ファクターとモメンタム（月次・日次）。
\item 経済政策不確実性指数（EPU、Baker・Bloom・Davis）と地政学リスク指数（GPR、Caldara・Iacoviello）の月次の値。
\item 論文に載っている数値例（Idzorek の 8 資産、Meucci の 6 か国、対象論文の表）。
\end{itemize}

論文の数値例は、論文の表と照らし合わせて再現できることを確かめてから使った。例えば Idzorek の例では、事後の期待リターンが論文の表と最大 \V{idz.erdiff}\% の差で一致する。

\section*{記号}

\begin{center}
\small
\begin{tabular}{p{2.2cm}p{11.4cm}}
\toprule
記号 & 意味 \\
\midrule
$N$、$K$ & 資産の数、ビューの数 \\
$r$（または $X$） & 資産のリターン（無リスク金利を引いた超過リターン）。$N$ 次元 \\
$\mu$ & 期待リターン（分からない量）。$N$ 次元 \\
$\Sigma$ & リターンの共分散行列（$N\times N$） \\
$\wmkt$ & 時価総額の比率（市場ポートフォリオ） \\
$\delta$（$\lambda$） & リスク回避度。論文によって $\delta$ と $\lambda$ の両方が使われる \\
$\Pi$（$\pi$） & 均衡リターン $\Pi=\delta\Sigma\wmkt$ \\
$\tau$ & 均衡の不確かさの大きさ。事前分布 $\mu\sim\mathcal N(\Pi,\tau\Sigma)$ \\
$P$ & ビューの組み合わせを表す行列（$K\times N$）。$k$ 行目を $p_k$ と書く \\
$Q$（$v$、$q$） & ビューの値（$K$ 次元） \\
$\Omega$ & ビューの不確かさ（$K\times K$、多くは対角）。対角成分を $\omega_k$ と書く \\
$\muBL$、$M$ & 事後の期待リターンと、その不確かさ（事後の共分散） \\
$C_k$ & Idzorek の確信度（0〜100\%） \\
$c$ & Meucci の確信度（$\Omega=\frac{1}{c}P\Sigma P\T$） \\
$s$ & ビューの採用率。均衡からビューへ何割動くか（1 資産なら $s=\tau\sigma^2/(\tau\sigma^2+\omega)$） \\
IC & 予想と実際のリターンの相関（情報係数） \\
TE、IR & トラッキングエラー（ベンチマークからのずれの標準偏差）、情報レシオ（超過リターン ÷ TE） \\
\bottomrule
\end{tabular}
\end{center}

\section*{要約：この資料で確かめたこと}

\begin{keybox}[要約]
\begin{enumerate}
\item BL の結果を決めるのは、$\Omega$ と $\tau\Sigma$ の大きさの比である。1 資産・1 ビューなら、均衡からビューへ動く割合は $s=\tau\sigma^2/(\tau\sigma^2+\omega)$ になる。$\Omega$ を $\tau$ に比例させると、事後の期待リターンから $\tau$ が消える（\chapref{ch:tau}）。
\item Idzorek の確信度の方法は、閉じた式 $\omega_k=\tau p_k\Sigma p_k\T(1-C_k)/C_k$ と同じ結果になる。ただし、指定した確信度どおりに比率が動くのはビューが 1 つのときだけで、3 つ同時に入れた例では、65\% と指定したビューが \V{idz.ach3}\% しか動かなかった（\chapref{ch:idzorek}）。
\item 予想の当たり具合を IC で測れば、$\Omega$ を決められる。Meucci の市場版の BL で $\omega=\sigma^2(1/\IC^2-1)$ と置くと、Grinold の式 $\alpha=\IC\cdot\sigma\cdot z$ と同じになる。IC が 0.05 なら、採用率は 0.25\% である（\chapref{ch:accuracy}）。
\item ビューが多いときは、$\Omega$ を対角にしてはいけない。誤差の相関を無視すると、銘柄の数だけ確信が積み上がる。本プロジェクトのデータでは、1 銘柄あたり数 \% のはずの採用率が、対角の $\Omega$ では中央値 \V{bt.calib.adopt}\% になった（\chapref{ch:project}）。
\item LLM の回答のばらつきは、当たり具合を表さない。回答の分散は、実際の誤差の 2 乗の中央値の \V{llm.ratio.med} 倍で、予想を 8 倍前後信じすぎる設定になる（\chapref{ch:project}）。
\item 政治・政策の指数（EPU、GPR）は、直近のボラティリティや VIX を入れると、翌月のボラティリティを説明する力がほとんど残らなかった。$\Omega$ を伸び縮みさせる材料としては、市場のボラティリティのほうが強い（\chapref{ch:politics}）。
\item モメンタムは、ボラティリティで規模を調整すると、急落が小さくなり、シャープレシオが \V{mom.sh.all} から \V{mom.sh.all.sc} に上がった。ビューの確信度を「モメンタム自身の最近のボラティリティ」で割るのは理にかなう（\chapref{ch:momentum}）。
\item 平均のビューだけなら、エントロピー・プーリングは市場版の BL とほぼ同じ答えを出す（差は最大 \V{ep.maxdiff}\%）。違いが出るのは、順位・ボラティリティ・下落確率のような、平均以外のビューを入れたときである（\chapref{ch:ep}）。
\item 予想が正しく縮められていても、平均分散最適化はトラッキングエラーを大きく取りがちである。トラッキングエラーの上限をはっきり置くほうがよい（\chapref{ch:te}、\chapref{ch:project}）。
\end{enumerate}
\end{keybox}
